% !TEX program = tectonic
\documentclass[9pt,a4paper]{extarticle}
\usepackage[top=0.8cm,bottom=0.8cm,left=0.9cm,right=0.9cm]{geometry}
\usepackage[T1]{fontenc}\usepackage[utf8]{inputenc}\usepackage{lmodern}
\usepackage[french]{babel}
\usepackage{amsmath,amssymb}\usepackage{textcomp}\usepackage{xcolor}\usepackage{booktabs}\usepackage{array}
\newcolumntype{R}[1]{>{\raggedright\arraybackslash}p{#1}}
\usepackage{enumitem}\usepackage{multicol}\usepackage{graphicx}
\usepackage[hidelinks]{hyperref}\usepackage{url}
\urlstyle{sf}
\graphicspath{{figs_nb11/}}
\definecolor{bleu}{RGB}{20,77,144}\definecolor{vert}{RGB}{20,110,60}
\definecolor{rouge}{RGB}{160,30,30}\definecolor{grisf}{RGB}{110,110,110}
\newcommand{\fig}[2]{\begin{center}\includegraphics[width=\linewidth]{#1}\end{center}\vspace{-2.5mm}{\scriptsize\color{grisf} #2}\par\vspace{0.3mm}}
\pagestyle{empty}\setlength{\parindent}{0pt}\setlength{\parskip}{1.1pt}
\setlength{\abovedisplayskip}{2.5pt}\setlength{\belowdisplayskip}{2.5pt}
\setlength{\abovedisplayshortskip}{1pt}\setlength{\belowdisplayshortskip}{1pt}
\renewcommand{\arraystretch}{0.95}
\setlength{\columnsep}{6mm}\setlength{\columnseprule}{0.3pt}
\renewcommand{\columnseprulecolor}{\color{grisf!40}}
\newcommand{\et}[2]{\vspace{1.6mm}{\color{bleu}\bfseries #1 $\cdot$ #2}\par\vspace{0.2mm}}
\newcommand{\obs}[1]{\vspace{0.2mm}\hangindent=3.4mm\hangafter=0{\color{bleu}\footnotesize\textbullet}~{\small #1}\par\vspace{0.2mm}}
\newcommand{\lire}[1]{{\scriptsize\color{vert}$\blacktriangleright$~\textbf{ce que ça dit :} #1}\par}
\newenvironment{pts}{\begin{itemize}[leftmargin=3.4mm,itemsep=0.2mm,topsep=0.4mm,parsep=0pt,label={\color{bleu}\footnotesize\textbullet}]}{\end{itemize}}
% ── renvois cliquables : ils ouvrent le PDF archivé (dossier docs_nanc_gop/) À LA PAGE exacte ──
\newcommand{\rf}[3]{\href{docs_nanc_gop/#2\#page=#3}{\color{bleu}\textsuperscript{\,[#1\,p.#3]}}}
\newcommand{\sPro}[1]{\rf{1}{1-prospectus-statutaire-SAI-20260127.pdf}{#1}}
\newcommand{\sOri}[1]{\rf{2}{2-prospectus-origine-20230203.pdf}{#1}}
\newcommand{\sChg}[1]{\rf{3}{3-changement-methode-20240802.pdf}{#1}}
\newcommand{\sAtrois}[1]{\rf{4}{4-rapport-annuel-FY2023.pdf}{#1}}
\newcommand{\sAquatre}[1]{\rf{5}{5-rapport-annuel-FY2024.pdf}{#1}}
\newcommand{\sAcinq}[1]{\rf{6}{6-rapport-annuel-FY2025.pdf}{#1}}
\newcommand{\sTick}[1]{\rf{7}{7-changement-code-KRUZ-GOP-20250318.pdf}{#1}}
\newcommand{\sSite}[1]{\rf{8}{8-fact-sheet-NANC-juin-2026.pdf}{#1}}
\newcommand{\sCoop}[1]{\rf{9}{9-cooper-la-regle-du-gerant-20231119.pdf}{#1}}
\newcommand{\sVen}[1]{\rf{10}{10-podcast-Venuto-mecanique-20241206.pdf}{#1}}
\begin{document}

\begin{center}
{\large\bfseries Répliquer NANC et GOP --- les ETF qui copient le Congrès}\\[0.4mm]
{\footnotesize\color{grisf} peut-on reconstruire ces portefeuilles à partir des données publiques ? y a-t-il un avantage à les suivre ?}\\[0.3mm]
{\scriptsize\color{grisf} Chaque affirmation sur le fonds porte un renvoi {\color{bleu}\textsuperscript{[n]}} \textbf{cliquable} vers le document officiel qui l'établit --- liste complète en fin de fiche.}
\end{center}
\vspace{-0.5mm}

\begin{multicols}{2}

% ═══════════ 1 ═══════════
\et{1}{La question}
\begin{pts}
\item Les élus du Congrès américain sont \textbf{obligés de déclarer} leurs transactions boursières (loi STOCK, 2012). Deux ETF vendent cette information au public : \textbf{NANC} copie les élus \textbf{démocrates}, \textbf{GOP} les \textbf{républicains}.
\item Ils sont nés le même jour (\textbf{06/02/2023})\sPro{13}, coûtent 0,72--0,73 \%/an\sPro{5}, et pèsent \textbf{282 M\$} (NANC, 104 lignes) et \textbf{88 M\$} (GOP, 138 lignes)\sSite{1}.
\item \textbf{Ce qu'on cherche} : (i) peut-on \textbf{reconstruire} leur portefeuille à partir des seules déclarations publiques ? (ii) cette stratégie a-t-elle un \textbf{avantage} sur le marché ?
\item \textbf{Pourquoi ce terrain est bon} : un ETF doit \textbf{publier son portefeuille chaque jour}\sSite{2}. On peut donc comparer notre réplique à la vérité, ligne par ligne.
\end{pts}

% ═══════════ 2 ═══════════
\et{2}{Ce à quoi on a accès}
\begin{center}\scriptsize
\begin{tabular}{@{}R{20mm} R{28mm} R{25mm}@{}}
\toprule
source & contenu & volume \\
\midrule
déclarations (PTR) & chaque transaction : titre, sens, date, \textbf{fourchette} de montant & \textbf{134\,464} transactions, 372 élus, 2014--2026 \\
holdings des ETF & leur portefeuille réel, publié chaque jour & 104 (NANC) et 139 (GOP) lignes \\
documents SEC & prospectus, SAI, rapports annuels, lettres du gérant & 2023 à 2026 \\
\bottomrule
\end{tabular}
\end{center}
\obs{\textbf{La limite décisive de nos données} : un élu ne déclare pas un montant, mais une \emph{fourchette}. Et \textbf{77 \% des achats tombent dans la même} (\og{}1\,001 -- 15\,000 \$\fg{}) : on leur attribue à tous 8\,000 \$. Pour les trois quarts du flux, \textbf{on ne connaît donc pas les montants réels}, seulement leur ordre de grandeur.}

% ═══════════ 3 ═══════════
\et{3}{Ce que le fonds dit faire --- et ce qu'il fait vraiment}
\textbf{Ce qu'annonce le prospectus} :
\begin{pts}
\item Univers : les titres achetés par les élus du parti, sur les \textbf{3 dernières années}, via les déclarations.\sPro{7}
\item Pondération au \og{}\textbf{niveau de trading}\fg{} : gros achats, achats \emph{répétés}, achats \emph{par plusieurs élus} $\to$ surpondérés ; petits, ventes récentes, coups isolés $\to$ sous-pondérés.\sPro{7}
\item Entre \textbf{100 et 200 lignes}. Aucun plafond de position n'est annoncé.\sPro{7}
\end{pts}
\textbf{Ce que révèlent les rapports annuels} --- et qu'on avait manqué en ne lisant que le résumé :
\begin{pts}
\item \textbf{Le fonds a changé de gérant et de méthode}, et c'est \emph{écrit} : le \textbf{02/08/2024} Tidal reprend la gestion\sChg{3}, le texte passe de \og{}\emph{créer le portefeuille \textbf{initial}}\fg{}\sOri{3} à \og{}\emph{\textbf{gérer} le portefeuille}\fg{}\sPro{7}, et le gérant écrit \og{}\emph{eliminated 582 very small positions}\fg{}\sAquatre{2} --- de $\sim$600 lignes à $\sim$170.
\item \textbf{Sa rotation le confirme} : elle s'effondre de \textbf{62 \% en 2024}\sAquatre{39} \textbf{à 10 \% en 2025}\sAcinq{27}. Il gérait activement, il ne touche presque plus à rien.
\item \textbf{Il y a du jugement humain} : les titres sont \og{}\emph{croisés avec les rôles en commission}\fg{} des élus\sAcinq{2}. Cela n'apparaît dans aucune règle publiée\sPro{7}.
\item \textbf{Le prospectus est incomplet} : aucun plafond n'y figure\sPro{7}, mais le fonds a \textbf{vendu 15 \% de son NVDA} en un an\sAcinq{17} --- il écrête en pratique.
\end{pts}
\begin{center}\scriptsize
\begin{tabular}{@{}l r r@{}}
\toprule
ce qu'ils publient & NANC & GOP \\
\midrule
performance depuis création, /an\sPro{13} & $+23{,}63\,\%$ & $+14{,}75\,\%$ \\
S\&P 500 (référence) & \multicolumn{2}{c}{$+20{,}96\,\%$} \\
\midrule
rotation FY2023\sAtrois{74} \emph{(7,8 mois)} & 44 \% & 46 \% \\
\quad \emph{soit, annualisée} & \emph{$\sim$68 \%} & \emph{$\sim$71 \%} \\
\quad FY2024\sAquatre{39} & 62 \% & 54 \% \\
\quad FY2025\sAcinq{27} & \textbf{10 \%} & \textbf{16 \%} \\
\bottomrule
\end{tabular}
\end{center}
\lire{NANC bat le S\&P de \textbf{2,7 points} par an ; \textbf{GOP le sous-performe de 6,2 points}. Ce ne sont pas deux réussites symétriques. La rotation FY2023 couvre 7,8 mois et n'est \textbf{pas annualisée} par le rapport\sAtrois{74} : la vraie séquence est $\sim$68 $\to$ 62 $\to$ 10 \%.}
\obs{\emph{L'objection} : l'actif moyen est au dénominateur, donc la croissance de l'encours pourrait à elle seule écraser le ratio. Les chiffres disent l'inverse\sAcinq{27} :}
\begin{center}\scriptsize
\begin{tabular}{@{}l r r r@{}}
\toprule
NANC & FY2023 & FY2024 & FY2025 \\
\midrule
encours (M\$) & 9,4 & 175,7 & 255,2 \\
croissance & --- & $\times$\textbf{18} & $+45\,\%$ \\
rotation & $\sim$68 \% & \textbf{62 \%} & \textbf{10 \%} \\
\bottomrule
\end{tabular}
\end{center}
\lire{la plus forte croissance ($\times$18) coïncide avec la rotation \textbf{la plus haute}, la plus faible avec la \textbf{plus basse} : l'artefact produirait exactement le contraire. \emph{Limite subsistante} : le chiffre \og{}exclut les transactions en nature\fg{}\sAcinq{27} --- 10 \% est un \textbf{plancher}.}

% ═══════════ 4 ═══════════
\et{4}{Ce qu'on fait --- trois lectures de leur règle}
\obs{\textbf{La règle, dite par le gérant lui-même.} Le prospectus reste vague (\og{}niveau de trading\fg{})\sPro{7}, mais \textbf{Christian Cooper}, co-gérant, l'a écrite en clair\sCoop{1} : \og{}\emph{build the holdings around the \textbf{relative dollars disclosed} by each member and their families. In other words, \textbf{more dollars equals more \og{}votes\fg{}}}\fg{} ; \og{}\emph{we explicitly choose the \textbf{midpoint} of that disclosure range}\fg{} ; \og{}\emph{this choice is what drives \textbf{all the relative allocations}}\fg{}. Soit : \textbf{poids $\propto$ dollars déclarés}, milieu de fourchette, familles incluses.}
On en tire trois lectures. Notation : $A_k$ le montant déclaré, $\delta_k$ la date de divulgation, $P_i$ le cours, $w_i$ un poids, $q_i$ un nombre de parts.
\begin{center}\scriptsize
\begin{tabular}{@{}R{16mm} R{26mm} R{26mm}@{}}
\toprule
 & \textbf{A --- le score} & \textbf{B --- le livre} \\
\midrule
principe & on \textbf{reclasse} tout chaque mois & on \textbf{achète à chaque déclaration}, puis on ne touche plus \\
ventes & ignorées & retirent des parts \\
sélection & les 150 plus gros scores & tout titre déclaré \\
plafond & 8,5 \% & aucun \\
poids & ramenés à la cible & \textbf{dérivent} avec les cours \\
\midrule
sa rotation & 66 \%/an & 13 \%/an \\
\textbf{quel régime} & \textbf{le fonds en 2023-24}\sAtrois{74}\sAquatre{39} (44--62 \%) & \textbf{le fonds en 2025}\sAcinq{27} (10 \%) \\
\bottomrule
\end{tabular}
\end{center}
\[ \text{A}:\quad s_i=\underbrace{B_i}_{\$\ \text{acheté}}\times\underbrace{m_i}_{\text{nb d'élus}},\qquad w_i\propto s_i\ \ (\text{plafonné}) \]
\[ \text{B}:\quad q_i(t)=q_i(t_0)+\!\!\sum_{\text{déclarations}}\!\!\pm\frac{A_k}{P_i(\delta_k)},\qquad w_i=\frac{q_iP_i}{\sum_j q_jP_j} \]
\obs{\textbf{C --- la bascule} : \textbf{A jusqu'à $t_1=$ 02/08/2024, puis B}, le portefeuille étant converti en parts à l'échelle du régime : $q_i(t_1)=w^{A}_i\,S/P_i(t_1)$, $S$ la taille du livre en \$. La date n'est pas choisie : c'est celle du changement de gérant (\S3). \textbf{C est notre version de référence.}}

% ═══════════ 5 ═══════════
\et{5}{Les résultats}
\textbf{(a) Reproduit-on leur portefeuille ?} Deux mesures, avec $w^{\star}_i$ leur poids \textbf{publié} et $T_{10}$ leurs 10 plus grosses lignes :
\[ \mathrm{expo}=\sum_i\min\!\bigl(w_i,\,w^{\star}_i\bigr),\qquad
   \mathrm{recouv}_{10}=\frac{\sum_{i\in T_{10}}\min\!\bigl(w_i,w^{\star}_i\bigr)}{\sum_{i\in T_{10}}w^{\star}_i} \]
\obs{\textbf{Pourquoi le $\min$} : il est \textbf{non compensatoire}. Une somme laisserait une surpondération racheter une sous-pondération --- mettre 8,5 sur GOOG (réel 5,8) \og{}compenserait\fg{} les 0,2 sur AMAT (réel 5,2). Avec $\min$, surpondérer ne rapporte rien et sous-pondérer coûte plein tarif. \emph{$\mathrm{expo}\in[0,1]$ : la fraction des deux portefeuilles qui est littéralement la même chose.}}
\begin{center}\scriptsize
\begin{tabular}{@{}l r r r r@{}}
\toprule
 & \multicolumn{2}{c}{NANC} & \multicolumn{2}{c}{GOP} \\
 & expo. & top-10 & expo. & top-10 \\
\midrule
A --- le score & \textbf{51 \%} & \textbf{65 \%} & \textbf{31 \%} & \textbf{25 \%} \\
B --- le livre & 46 \% & 51 \% & 26 \% & 14 \% \\
C --- la bascule & 48 \% & 57 \% & 29 \% & 20 \% \\
\bottomrule
\end{tabular}
\end{center}
\lire{on reproduit \textbf{la moitié} du portefeuille NANC, \textbf{moins d'un tiers} du GOP. \emph{Contrôle} : la corrélation des poids donne le même classement ($\rho=0{,}75$ NANC, $0{,}26$ GOP) --- mais sur les \textbf{rangs} elle tombe à $0{,}06$, on a leurs gros noms sans leur ordre.}
\obs{\textbf{Ce n'est pas un coup de chance sur une date.} Les états \emph{N-PORT} déposés à la SEC donnent leur portefeuille réel à \textbf{13 dates} (03/2023 $\to$ 03/2026). Sur ces 13 photos, l'exposition commune reste \textbf{40--53 \%} pour NANC et \textbf{31--39 \%} pour GOP --- stable avant comme après le reset d'août 2024.}
\obs{\textbf{Et la règle de Cooper ne fait pas mieux que la nôtre.} Sur ces mêmes 13 dates : \emph{dollars purs} (sa formulation) contre \emph{dollars $\times$ nombre d'élus} (notre score) $\to$ \textbf{44,5 \% vs 46,8 \%} (NANC) et \textbf{35,2 \% vs 34,0 \%} (GOP). \textbf{1 à 2 points d'écart} : le choix de pondération n'est \emph{pas} ce qui nous sépare d'eux --- ce qui renforce le \S6.}

\textbf{(b) Où est-ce que ça passe, où est-ce que ça casse ?} Leurs 10 plus grosses lignes, en \% :
\begin{center}\scriptsize
\begin{tabular}{@{}l r r r r@{}}
\toprule
NANC & eux & A & B & C \\
\midrule
NVDA & 8,53 & 8,50 & 8,93 & 11,73 \\
GOOG & 5,81 & 8,50 & 4,96 & 8,94 \\
MSFT & 5,29 & 8,22 & 3,81 & 5,55 \\
\textbf{AMAT} & 5,19 & \textbf{0,21} & 0,31 & 0,42 \\
AMZN & 4,70 & 8,24 & 3,05 & 2,88 \\
AAPL & 4,58 & 6,57 & 0,61 & 2,38 \\
\textbf{CRWD} & 3,41 & \textbf{0,13} & 1,38 & 0,07 \\
\textbf{PM} & 3,20 & \textbf{0,11} & 0,05 & 0,04 \\
AXP & 2,84 & 0,15 & 0,65 & 0,79 \\
LLY & 2,71 & 0,75 & 0,43 & 0,36 \\
\bottomrule
\end{tabular}
\end{center}
\lire{on retrouve \textbf{leurs grands noms} (NVDA, GOOG, MSFT, AMZN, AAPL) mais on se trompe sur \textbf{les tailles} ; et on \textbf{rate complètement} AMAT, CRWD, PM. Côté GOP c'est pire : leur \textbf{première position, FIX (8,4 \%), on la met à 0}.}
\fig{nanc_gop_v3.png}{leur portefeuille (gris) contre le nôtre : les grands noms coïncident à gauche (NANC), presque rien ne coïncide à droite (GOP).}

\textbf{(c) La rotation --- ce qui rattache chaque mécanique à sa période.} Le rapport annuel la définit ainsi\sAtrois{74} : \og{}\emph{le numérateur inclut le plus petit des achats ou des ventes ; le dénominateur, la valeur moyenne des positions}\fg{} --- soit
\[ \tau \;=\; \frac{\min\bigl(\textstyle\sum\text{achats},\ \sum\text{ventes}\bigr)}{\text{actif moyen}} \]
Chez nous les portefeuilles sont en \textbf{poids}. Entre deux coupes, les poids bougent pour deux raisons : les cours \emph{dérivent} (sans aucune transaction), puis on \emph{rebalance} (là il y a transaction). On isole donc la seconde :
\[ \begin{gathered}
w^{\mathrm{d}}_i(t)=\frac{w_i(t^-)\,P_i(t)/P_i(t^-)}{\sum_j w_j(t^-)\,P_j(t)/P_j(t^-)} \\[0.6mm]
\tau_Y=\sum_{m\in Y}\tfrac12\sum_i\bigl|w^{\text{cible}}_i(m)-w^{\mathrm{d}}_i(m)\bigr|
\end{gathered} \]
\obs{\textbf{Pourquoi le $\tfrac12$} : $\sum_i w_i=1$ avant comme après, donc ce qu'on achète égale ce qu'on vend et $\sum_i|\Delta w_i|$ compte le mouvement \textbf{deux fois} --- la moitié restitue $\min(\text{achats},\text{ventes})$. \textbf{Pourquoi pas de dénominateur} : un poids \emph{est} déjà une fraction de l'actif. Les deux définitions coïncident donc.}
\begin{center}\scriptsize
\begin{tabular}{@{}l r r r@{}}
\toprule
rotation/an & 2023 & 2024 & 2025 \\
\midrule
A --- le score & 59 \% & 67 \% & 82 \% \\
B --- le livre & 13 \% & 11 \% & 16 \% \\
C --- la bascule & 48 \% & 37 \% & 26 \% \\
\textbf{le fonds} & \textbf{$\sim$68 \%} & \textbf{62 \%} & \textbf{10 \%} \\
\bottomrule
\end{tabular}
\end{center}
\lire{\textbf{A} est compatible avec 2023-24 et incompatible avec 2025 ; \textbf{B}, l'inverse --- d'où \textbf{C}. Un scalaire par an \textbf{exclut} des mécaniques, il n'en \textbf{identifie} aucune : le changement de régime est établi au \S3, la rotation le date. \textbf{Limite de C} : elle a le \emph{sens} de la chute, pas l'\emph{ampleur} (26 \% contre 10 \%).}

\textbf{(d) Y a-t-il un avantage ?} Pour l'année $Y$, puis moyenne sur les années ; $\beta$ par régression sur le marché en excès :
\[ \begin{gathered}
x_Y=\prod_{t\in Y}(1{+}r_t)-\prod_{t\in Y}(1{+}r^{\mathrm{SPY}}_t),\qquad t=\frac{\bar x}{\sigma_x/\sqrt{n}} \\[0.6mm]
r-r_f=\alpha+\beta\bigl(r^{\mathrm{SPY}}-r_f\bigr)+\varepsilon
\end{gathered} \]
\begin{center}\scriptsize
\begin{tabular}{@{}l l r r r@{}}
\toprule
 & période & excès/an & $t$ & $\beta$ \\
\midrule
A NANC & 2014-26 & $+2{,}98\,\%$ & 1,47 & 1,08 \\
A GOP & 2014-26 & $+1{,}24\,\%$ & 1,62 & 1,00 \\
C NANC & 2023-26 & $+2{,}46\,\%$ & 0,98 & 1,06 \\
C GOP & 2023-26 & $-0{,}27\,\%$ & $-0{,}15$ & 0,93 \\
\bottomrule
\end{tabular}
\end{center}
\lire{$\beta\approx1$ partout et $|t|<2$ partout : ces portefeuilles \textbf{sont} le marché, et \textbf{aucun excès n'est distinguable du hasard}.}

% ═══════════ 6 ═══════════
\et{6}{Pourquoi ça ne colle pas}
On compare enfin des \textbf{dollars} et non des pourcentages. Soit $F_i$ le \textbf{\$ total d'achats jamais déclaré} par le parti sur le titre $i$ :
\[ \lambda_i \;=\; \frac{w^{\star}_i\times\mathrm{AUM}}{F_i}
   \qquad\text{poids}\propto\text{flux}\iff \lambda_i\ \text{constant} \]
\begin{center}\scriptsize
\begin{tabular}{@{}l r r r@{}}
\toprule
 & leur position & flux déclaré & rapport \\
\midrule
NVDA & 24,1 M\$ & 23,0 M\$ & \textbf{1,0} \\
GOOG & 16,4 M\$ & 16,4 M\$ & \textbf{1,0} \\
AMZN & 13,3 M\$ & 13,8 M\$ & \textbf{1,0} \\
MSFT & 14,9 M\$ & 83,0 M\$ & 0,2 \\
\textbf{AMAT} & 14,6 M\$ & 2,2 M\$ & \textbf{6,7} \\
\textbf{PM} & 9,0 M\$ & 1,0 M\$ & \textbf{9,0} \\
\textbf{FIX} \emph{(GOP)} & 7,5 M\$ & 0,3 M\$ & \textbf{25,0} \\
\bottomrule
\end{tabular}
\end{center}
\lire{$\lambda$ va de \textbf{0,2 à 25} au lieu d'être constant --- et les seuls titres où $\lambda\approx1$ sont \emph{exactement} ceux qu'on reproduit bien. Le fonds détient \textbf{25 fois} plus de FIX que tout le flux jamais déclaré dessus : aucune règle fondée sur les déclarations publiques ne peut produire ça.}
\obs{\textbf{Correction importante : leur donnée est la nôtre.} On a d'abord cru à une source privilégiée. C'est faux, et le gérant le dit\sVen{1} : \og{}\emph{the data is \textbf{public disclosure}. \textbf{Anybody can go get it}}\fg{} --- Unusual Whales ne fait que la rendre exploitable. L'écart ne vient donc \textbf{pas} des données d'entrée, mais de ce qu'ils en font : \og{}\emph{\textbf{Dan Weiskopf spends a lot of time poking through the data to try and find the patterns}}\fg{} (auditions en commission, achats simultanés). \textbf{C'est un jugement humain, pas une formule} --- et c'est ce qui produit la dispersion $\times$25 ci-dessus.}
\obs{\textbf{Et ils écartent un bruit que nous comptions.}\sVen{1} \og{}\emph{Some of them have a broker who's just doing \textbf{direct indexing}\ldots{} when we see \textbf{300 trades and they were all, like, 5 shares of something, we know that's just a rebalance}}\fg{}. Or ces dépôts en masse pèsent \textbf{66 \% du \$ démocrate} et \textbf{85 \% du républicain} chez nous. En les retirant (seuil de l'ordre de 50--100 opérations le même jour, cohérent avec leurs \og{}300 trades\fg{}), la reproduction du \textbf{top-10} progresse nettement : \textbf{43,5 $\to$ 56,6 \%} pour NANC et \textbf{25,9 $\to$ 33,9 \%} pour GOP (exposition GOP $+5{,}7$ pt). \emph{Filtrer plus fin dégrade} --- et calibrer le seuil sur la cible serait de l'ajustement, pas une réplication.}
\obs{\textbf{Contre-épreuve --- ce n'est pas notre faute.} On a audité et réparé notre propre chaîne : \textbf{7\,287 transactions} écartées par un filtre ont été réintégrées, \textbf{10 séries de prix} manquantes récupérées, les montants \og{}$>$1 M\$\fg{} revalorisés. \textbf{La fidélité ne bouge pas.} On a aussi testé l'hypothèse \og{}ils utilisent les déclarations de \emph{patrimoine} et non de transactions\fg{} : \textbf{réfutée} --- AMAT y pèse 0,009 \% contre 5,19 \% du fonds.}

% ═══════════ 7 ═══════════
\et{7}{Conclusion}
\begin{pts}
\item \textbf{Ces produits ne sont pas reproductibles à partir des données publiques.} On retrouve leurs grands noms et l'ordre de grandeur de leur performance, mais pas leur portefeuille : la moitié pour NANC, moins d'un tiers pour GOP.
\item \textbf{On sait pourquoi, et c'est chiffré} : (i) leurs positions valent jusqu'à \textbf{25 fois} le flux déclaré correspondant, alors même que \textbf{la donnée de départ est la même que la nôtre}\sVen{1} : l'écart est produit par leur \textbf{sélection}, pas par leurs entrées ; (ii) le gérant ajoute un \textbf{jugement discrétionnaire} qu'il admet lui-même\sAcinq{2} ; (iii) le fonds a \textbf{changé de mécanique} en août 2024\sChg{3}\sAquatre{2}.
\item \emph{Appui extérieur} : Baulkaran \& Jain, \emph{Economics Letters} 250 (2025), concluent indépendamment que \og{}\emph{neither significantly outperforms market returns}\fg{}.
\item \textbf{Et surtout : il n'y a pas d'avantage.} $\beta\approx1$, aucun $t$ au-dessus de 2. Le fonds démocrate bat le S\&P de 2,7 points par an, le républicain le sous-performe de 6,2\sPro{13} : c'est une \textbf{exposition au marché} avec un biais sectoriel, pas une compétence.
\end{pts}

\end{multicols}

% ═══════════ ANNEXE ═══════════
\vspace{1.5mm}\hrule\vspace{1.5mm}
{\color{bleu}\bfseries\large Annexe --- la règle, dans leurs mots}\par\vspace{1.5mm}
\textbf{Prospectus}\sPro{7} \emph{(\og{}Unusual Whales Subversive Democratic Trading ETF\fg{}, \og{}Principal Investment Strategies\fg{}, p.\,2 ; la version républicaine est identique, \og{}Republican\fg{} au lieu de \og{}Democratic\fg{})} :
{\footnotesize\itshape
``To manage the Fund's portfolio, the Adviser will obtain and use information derived by the Data Provider from PTRs filed by Democratic U.S. Congresspeople for the past 3 years. Purchases made during that time will be netted against any sales of the same security when managing the Fund's portfolio of equity securities. As the investment thesis of the Fund is to track the trading activity of Democratic U.S. Congresspeople while in office, equity securities acquired by Democratic U.S. Congresspeople prior to his or her swearing in (or the 3-year lookback period) are not considered when managing the Fund's portfolio. To the extent a Democratic U.S. Congressperson sells equity securities that were acquired prior to his or her swearing in, the Adviser will not adjust the Fund's portfolio.

Under normal circumstances, the Fund will invest in a portfolio of between 100 to 200 holdings. However, the number and size of positions held by the Fund will vary based on the number of positions traded by Democratic U.S. Congresspeople. Individual holdings will be weighted based on the level of reported trading in a security by Democratic U.S. Congresspeople. Securities with large purchases, recurring purchases and purchases from multiple Democratic U.S. Congresspeople will be overweighted. Securities with small purchases, recent sales and one-off trades by Democratic U.S. Congresspeople will be excluded or underweighted. The Adviser may exclude positions that are traded by Democratic U.S. Congresspeople at its discretion. Considerations for excluding a security include, but are not limited to, limited liquidity of such security and pending corporate actions that may impact such security.''\par}
\vspace{1.5mm}
\textbf{Lettre du gérant}\sAquatre{2} \emph{(rapport annuel N-CSR, exercice clos le 30/09/2024)} --- le reset :
{\footnotesize\itshape ``eliminated 582 very small positions''\par}
\vspace{1mm}
\textbf{Lettre du gérant}\sAcinq{2} \emph{(rapport annuel N-CSR, exercice clos le 30/09/2025)} --- la part discrétionnaire :
{\footnotesize\itshape ``cross referenced against the Democratic members' committee roles''\par}
\vspace{1mm}
\textbf{Prospectus, section \og{}Data Provider\fg{}}\sPro{34} --- qui construit réellement le portefeuille :
{\footnotesize\itshape ``UW plays no role in the construction or implementation''\par}

% ═══════════ SOURCES ═══════════
\vspace{2mm}\hrule\vspace{1.5mm}
{\color{bleu}\bfseries\large Sources --- les documents eux-mêmes, à la page exacte}\par\vspace{1mm}
{\scriptsize\color{grisf} Les renvois {\color{bleu}[n\,p.X]} du texte \textbf{ouvrent directement le PDF} archivé dans \path{docs_nanc_gop/}, à la page indiquée. Ces PDF sont les documents officiels tels que déposés, convertis sans retouche ; l'adresse d'origine est donnée pour chacun. Les fonds relèvent de \emph{Tidal Trust I} (CIK 1742912) depuis fin 2024 et de \emph{Series Portfolios Trust} (CIK 1650149) auparavant : c'est pourquoi les dépôts sont répartis sur deux identifiants. Liens vérifiés le 27/07/2026.\par}
\vspace{1mm}

{\footnotesize
\textbf{[1] Prospectus statutaire $+$ SAI} --- 485BPOS, 27/01/2026, 146 p. \emph{Ce qu'on y vérifie} : la règle complète (\textbf{p.\,7} pour NANC, p.\,17 pour GOP), les frais (p.\,5), la performance depuis création (\textbf{p.\,13}), et le rôle du fournisseur de données (\textbf{p.\,34}).\\
\path{docs_nanc_gop/1-prospectus-statutaire-SAI-20260127.pdf} \quad{\color{grisf}origine :}\\ \url{https://www.sec.gov/Archives/edgar/data/1742912/000199937126001797/nanckruz-485bpos_012726.htm}

\vspace{1mm}
\textbf{[2] Prospectus d'origine} --- 497K, 03/02/2023, 9 p. \emph{Ce qu'on y vérifie} : la version qui dit \og{}\emph{to create the Fund's \textbf{initial} portfolio}\fg{} (\textbf{p.\,3}) --- la fenêtre 3 ans ne servait alors qu'au portefeuille de départ.\\
\path{docs_nanc_gop/2-prospectus-origine-20230203.pdf} \quad{\color{grisf}origine :}\\ \url{https://www.sec.gov/Archives/edgar/data/1650149/000089418923000948/unusualwhalessubversivedem.htm}

\vspace{1mm}
\textbf{[3] Changement de méthode} --- 497, 02/08/2024, 10 p. \emph{Ce qu'on y vérifie} : Tidal devient gérant et le portefeuille passe à 100--200 lignes (\textbf{p.\,3}).\\
\path{docs_nanc_gop/3-changement-methode-20240802.pdf} \quad{\color{grisf}origine :}\\ \url{https://www.sec.gov/Archives/edgar/data/1650149/000089418924004591/subversivecapitalkruznanc-a.htm}

\vspace{1mm}
\textbf{[4] Rapport annuel FY2023} --- N-CSR, exercice clos le 30/09/2023, 97 p. \emph{Ce qu'on y vérifie} : rotation \textbf{44 \%} pour NANC (\textbf{p.\,74}) et \textbf{46 \%} pour GOP (p.\,76), dans les \emph{Financial Highlights}.\\
\path{docs_nanc_gop/4-rapport-annuel-FY2023.pdf} \quad{\color{grisf}origine :}\\ \url{https://www.sec.gov/Archives/edgar/data/1650149/000089853123000485/uwsctetf-ncsra.htm}

\vspace{1mm}
\textbf{[5] Rapport annuel FY2024} --- N-CSR, exercice clos le 30/09/2024, 59 p. \emph{Ce qu'on y vérifie} : la phrase du gérant \og{}\emph{eliminated 582 very small positions}\fg{} (\textbf{p.\,2}) ; rotation \textbf{62 \%} NANC (\textbf{p.\,39}) et 54 \% GOP (p.\,40).\\
\path{docs_nanc_gop/5-rapport-annuel-FY2024.pdf} \quad{\color{grisf}origine :}\\ \url{https://www.sec.gov/Archives/edgar/data/1650149/000183988224044132/whales-ncsr_93024.htm}

\vspace{1mm}
\textbf{[6] Rapport annuel FY2025} --- N-CSR, exercice clos le 30/09/2025, 46 p. \emph{Ce qu'on y vérifie} : \og{}\emph{cross referenced against the \ldots{} committee roles}\fg{} (\textbf{p.\,2}) ; rotation \textbf{10 \%} NANC (\textbf{p.\,27}) et 16 \% GOP (p.\,28) ; et le portefeuille détaillé où l'on lit \textbf{142\,778 parts de NVIDIA} (\textbf{p.\,17}), contre 167\,829 l'année précédente.\\
\path{docs_nanc_gop/6-rapport-annuel-FY2025.pdf} \quad{\color{grisf}origine :}\\ \url{https://www.sec.gov/Archives/edgar/data/1742912/000199937125019807/suw-ncsr_093025.htm}

\vspace{1mm}
\textbf{[7] Changement de code} --- 497, 18/03/2025, 1 p. \emph{Ce qu'on y vérifie} : le fonds républicain passe de \textbf{KRUZ à GOP} --- indispensable pour retrouver ses dépôts antérieurs.\\
\path{docs_nanc_gop/7-changement-code-KRUZ-GOP-20250318.pdf} \quad{\color{grisf}origine :}\\ \url{https://www.sec.gov/Archives/edgar/data/1742912/000199937125002790/kruz-497_031825.htm}

\vspace{1mm}
\textbf{[9] La règle, écrite par le gérant} --- Christian H. Cooper (co-gérant), Unusual Whales 19/11/2023 ; version longue dans \emph{Barron's} 18/12/2023. \emph{Ce qu'on y vérifie} (\textbf{p.\,1}) : \og{}\emph{more dollars equals more votes}\fg{}, le \textbf{midpoint} des fourchettes, et \og{}\emph{this choice is what drives all the relative allocations}\fg{}.\\
\path{docs_nanc_gop/9-cooper-la-regle-du-gerant-20231119.pdf} \quad{\color{grisf}origine :}\\ \url{https://unusual-whales.ghost.io/unusual_blog/post/nanc-is-beating-the-s-p500-per-subversive/}

\vspace{1mm}
\textbf{[10] La mécanique sous Tidal, par le gérant actuel} --- Michael Venuto, \emph{The Meb Faber Show} \#560, 06/12/2024 (chapitre à 42:01). \emph{Ce qu'on y vérifie} (\textbf{p.\,1}) : la donnée est publique ; la sélection est \textbf{discrétionnaire} ; et le fonds \textbf{écarte le bruit} des comptes gérés en \emph{direct indexing}.\\
\path{docs_nanc_gop/10-podcast-Venuto-mecanique-20241206.pdf} \quad{\color{grisf}origine :}\\ \url{https://www.themebfabershow.com/episodes/3395kHHKQ3v}

\vspace{1mm}
\textbf{[8] Fiches produit de l'émetteur} --- au 30/06/2026, 2 p. chacune. \emph{Ce qu'on y vérifie} : l'encours, le nombre de lignes et la performance affichée.\\
\path{docs_nanc_gop/8-fact-sheet-NANC-juin-2026.pdf} \ et \ \path{docs_nanc_gop/8-fact-sheet-GOP-juin-2026.pdf}\\
{\color{grisf}origine :} \url{https://subversiveetfs.com/nanc/} \ et \ \url{https://subversiveetfs.com/gop/}
\par}

\vspace{1.5mm}
{\footnotesize\textbf{Données utilisées de notre côté} --- \emph{pour la réplique et la comparaison} :
\begin{pts}
\item \textbf{Portefeuilles réels des deux ETF} au 24/07/2026, tels que publiés quotidiennement par l'administrateur (Tidal Financial Group) : \path{cache/NANC_holdings_20260724.csv} (104 lignes) et \path{cache/GOP_holdings_20260724.csv} (139 lignes).
\item \textbf{Déclarations des élus} (PTR), table nettoyée du projet : \path{00_S1S2_donnees/data/clean/transactions_backtest_2014_2026.csv} --- 134\,464 transactions, 372 élus, 2014--2026.
\item \textbf{Tous les calculs} de cette fiche sont produits par le notebook \path{11_Strategie_Ticker.ipynb} (sections §13.4 à §13.12) : chaque nombre affiché ici est une sortie de cellule, aucun n'est recopié à la main.
\end{pts}\par}

\end{document}
